\section{Results Discussion}
\label{sec:results}

\subsection{QR Factorization and Linear Least-Squares Solution}
\label{subsec:problem1-results}

The modified Gram--Schmidt algorithm was applied to the
overdetermined system defined in Section~\ref{sec:problem}. The
computed factors satisfy $A\approx QR$, where
$Q\in\mathbb{R}^{5\times4}$ has approximately orthonormal columns
and $R\in\mathbb{R}^{4\times4}$ is upper triangular.

The computed matrices, rounded to six decimal places, are

\begin{equation}
    Q\approx
    \begin{bmatrix}
        0.295599  & 0.242348  & -0.209363 & 0.159387  \\
        0.197066  & -0.200200 & 0.322900  & 0.898629  \\
        0.492665  & 0.140811  & 0.781674  & -0.352688 \\
        -0.098533 & 0.938737  & -0.000348 & 0.202929  \\
        -0.788263 & 0.011495  & 0.490803  & 0.038631
    \end{bmatrix},
    \label{eq:computed-Q}
\end{equation}

and

\begin{equation}
    R\approx
    \begin{bmatrix}
        10.148892 & -4.926646 & -3.251587 & -4.236916 \\
        0         & 10.135490 & 2.859317  & 4.846950  \\
        0         & 0         & 13.048045 & 2.480397  \\
        0         & 0         & 0         & 10.020142
    \end{bmatrix}.
    \label{eq:computed-R}
\end{equation}

The quality of the factorization was assessed using the relative
reconstruction error and the orthogonality error. The computed
values were

\begin{equation}
    \eta_{\mathrm{QR}}=1.0683\times10^{-16},
    \qquad
    \eta_Q=4.5180\times10^{-16}.
\end{equation}

Both indicators are close to the expected scale of floating-point
rounding errors for this computation. These results support the
numerical accuracy of the factorization and the orthogonality of
the computed columns of $Q$.

The least-squares solution was obtained by solving
$R\mathbf{x}^{*}=Q^T\mathbf{b}$ using back substitution. The
resulting solution, rounded to nine decimal places, is

\begin{equation}
    \mathbf{x}^{*}\approx
    \begin{bmatrix}
        -0.404938762 \\
        -0.040220457 \\
        0.992366790  \\
        -0.568506366
    \end{bmatrix}.
    \label{eq:computed-solution}
\end{equation}

The Euclidean norm of the residual was

\begin{equation}
    \|\mathbf{A}\mathbf{x}^{*}-\mathbf{b}\|_2
    \approx 7.677354.
\end{equation}

The nonzero residual indicates that the original system is
inconsistent. This is expected for an overdetermined least-squares
problem and does not, by itself, indicate a failure of the numerical
method. Instead, the solution minimizes the residual norm among
all vectors in $\mathbb{R}^4$.

To assess the optimality of the computed solution, the norm of
the residual projected onto the column directions of $A$ was also
evaluated. The result was

\begin{equation}
    \left\|A^T(A\mathbf{x}^{*}-\mathbf{b})\right\|_2
    \approx 6.1020\times10^{-14}.
\end{equation}

This value is numerically small and is consistent with the
least-squares optimality condition $A^T(A\mathbf{x}^{*}-\mathbf{b})
    =\mathbf{0}$. Hence, the computed residual is approximately
orthogonal to the column space of $A$.

Finally, the diagonal entries of $R$ were compared with the
scale-dependent numerical tolerance used by the implementation.
All four diagonal entries were substantially larger than this
tolerance, supporting the numerical assessment that $A$ has full
column rank. Therefore, the least-squares solution is unique.

These results validate the use of the modified Gram--Schmidt
factorization and back substitution for the matrix considered in
this experiment. Nevertheless, modified Gram--Schmidt without
column pivoting is not a universally reliable rank-revealing
method for severely ill-conditioned matrices. Householder QR
or rank-revealing factorizations may be preferable when stronger
numerical robustness is required
\cite{golub2013matrix,higham2002accuracy}.